\section{从 Architecture 到 Universal Substrate}

讲座最后把 Transformer 放回 consolidation thesis。它已从 translation backbone 扩展到 vision、contrastive learning、语言模型和 tool-using systems，但课程也明确列出未覆盖的 layer norm、FFN 修改、higher-order attention 与 faster decoding。Universal substrate 是一种趋势描述，不是“所有问题已经统一”的宣告。

\subsection{其他领域的扩展与本讲的有意省略}

Vision Transformer 把图像 patch 当 token，contrastive model 把跨模态表示对齐；这些变化证明 token mixing + feature transformation 的 block 可迁移。与此同时，pre/post layer norm、gated FFN、higher-order interaction、speculative/non-autoregressive decoding 都会显著影响现代系统，本讲只选择 attention/position/system 主线。

\lecturefigure{slide-37-advancements-other-areas.jpg}{Transformer 扩展到 vision 与 contrastive learning}{官方录像恢复幻灯片 V037；00:52:33}

\begin{knowledgebox}{读图：迁移的是 block motif，不是输入天然相同}
图像需要 patch/tokenization 和二维 position，contrastive learning 需要跨样本 objective。先看共享的 Transformer encoder，再看 modality-specific input/output；跨领域成功支持 architecture reuse，却不意味着所有 inductive bias 都应删除。
\end{knowledgebox}

\lecturefigure{slide-38-not-covered-topics.jpg}{本讲未展开但同样重要的 Transformer 演化方向}{官方录像恢复幻灯片 V038；00:54:15}

\begin{warningbox}{一场 80 分钟讲座不能替代完整领域地图}
Pre-vs-post norm 影响深层稳定性，FFN/gating 影响 capacity，higher-order attention 改变 interaction，speculative decoding 改善 inference。把“未覆盖”误读为“不重要”，会让讲义比课堂更武断。
\end{warningbox}

\teachervoice{讲者主动展示 Not Covered 页，是一种 evidence hygiene：他承认大量重要进展没有进入叙事。高质量笔记也应保留范围声明，而不是为了显得完整，把所有后续术语硬塞进课堂。}

\subsection{LLM as a Prompt-programmable Substrate}

当同一 pretrained model 能通过 prompt 执行广泛任务，它开始像一种通用计算 substrate：用户不再为每个任务训练独立模型，而是用 context 指定行为。这里的“programming”仍是 probabilistic conditioning，受到训练数据、instruction following、context window 和 reliability 限制。

\lecturefigure{slide-39-llms-universal-substrate.jpg}{LLM 作为 prompt-programmable universal substrate}{官方录像恢复幻灯片 V039；00:57:15}

\begin{knowledgebox}{读图：Prompt 是接口，不是可靠性证明}
Pac-Man 式图示表达一个模型吸收多种任务。先问任务是否真的共享 model weights，再问输出能否验证；few-shot breadth 说明 consolidation，却不能保证 exact computation、fresh knowledge、long-horizon execution 或 safety。
\end{knowledgebox}

\begin{importantbox}{Consolidation 会把问题向下游迁移}
模型接口变统一后，差异不会消失，而会转移到 data mixture、prompt/tool schema、evaluation、serving、feedback 与 safety。系统更简单的是表面 API，不一定是完整工程栈。
\end{importantbox}

\subsection{现代 Transformer 是 Data-center System}

大模型训练和 serving 跨成千上万 GPU，涉及 all-reduce、all-gather、reduce-scatter 等\term{collectives（集合通信）}、网络拓扑、拥塞、并行策略与容错。模型图上的一个 layer，在运行时被切成多个 device 上的参数、activation 和 communication schedule。

\lecturefigure{slide-40-modern-transformer-datacenter.jpg}{现代大规模 Transformer 必须作为 data-center system 理解}{官方录像恢复幻灯片 V040；00:59:06}

\begin{knowledgebox}{读图：参数规模扩大后，网络也进入模型 critical path}
机房照片不是装饰。先问 tensor/data/pipeline parallel 如何切分，再看 collective 是否被 compute overlap；模型 quality 可能相同，cluster topology 和 congestion 却决定训练是否可扩展。Architecture researcher 也需要理解 systems boundary。
\end{knowledgebox}

\teachervoice{讲者说如果真的研究 modern large-scale Transformer，就必须学习 all-reduce、InfiniBand、RoCE 与 congestion。大 Transformer “now in some sense is a data center”；这与开场对 Dartmouth machine-capacity 判断的反转形成闭环。}

\subsection{Tools：哪些能力应该在模型内，哪些应该调用外部系统}

\term{tool use（工具使用）}让模型通过结构化 interface 调用 calculator、search、database、code executor 或其他 service。它不是回到脆弱 pipeline，而是把可验证、精确、已有的软件能力留在外部，让模型负责选择、组合和解释。

\lecturefigure{slide-41-llms-using-tools.jpg}{LLM 通过工具连接搜索、计算与外部软件}{官方录像恢复幻灯片 V041；00:59:24}

\begin{knowledgebox}{读图：Swiss-army knife 仍需要明确接口}
不同工具有 schema、权限、cost 和 failure mode。先看模型何时决定调用，再看结果如何验证和进入下一步；tool breadth 可以扩大 capability，但错误 routing、prompt injection 与不可逆 action 也扩大风险。
\end{knowledgebox}

\teachervoice{Q\&A 中 Vaswani 给出非常工程化的判断：不应花“每个位置十亿 FLOPs”去计算两个数，某些能力就该交给外部工具；模型应承担需要语义判断和组合的 thinking。工具边界是 capability allocation，不是模型不够纯粹。}

\subsection{It Is Time to Build：产品反馈也是学习信号}

讲者主张与真实用户互动可以暴露能力缺口，形成 data/feedback/model improvement loop。产品并不自动产生好数据：反馈会有 selection bias、privacy、reward hacking 和短期 KPI；但离开真实 workflow，研究者很难知道模型在哪些 interface 上真正失败。

\lecturefigure{slide-42-time-to-build.jpg}{真实使用与模型改进可以形成产品反馈闭环}{官方录像恢复幻灯片 V042；01:00:06}

\begin{importantbox}{Build 的研究价值来自可观测 Failure}
高价值闭环应记录任务、工具轨迹、用户修正、执行结果和失败分类，并区分“用户喜欢”与“任务正确”。只有把反馈转成可审计训练/evaluation artifact，产品使用才会提升模型，而不是只放大已有偏差。
\end{importantbox}

\subsection{本章小结}

Transformer 的 consolidation 已跨越模态和任务，但仍有大量未覆盖的 architecture work。LLM 成为通用 substrate 后，系统问题转向 data center、tools、feedback 与 verification；统一模型没有让工程消失，而是重新组织工程边界。

